% 原试卷 PDF 第52～54页。
\setchoices{label=(\Alph*)}

\begin{question}[index=1]
曲线 $\displaystyle y=\frac{\mathrm e^x+1}{2\mathrm e^x-1}$
\begin{choices}[columns=2]
\item 没有渐近线．
\item 仅有1条渐近线．
\item 仅有2条渐近线．
\item 仅有3条渐近线．
\end{choices}
\end{question}
\QuestionBreak

\begin{question}[index=2]
若二元函数 $f(x,y)$ 与 $g(x,y)$ 在点 $P_0(x_0,y_0)$ 的某邻域内有定义，$f(x,y)$ 在点 $P_0$ 处可微且 $f(x_0,y_0)=0$，$g(x,y)$ 在点 $P_0$ 处连续但不可微，则对于下列4个函数：
\[
\begin{gathered}
F_1(x,y)=f(x,y)|f(x,y)|,\quad F_2(x,y)=f(x,y)g(x,y),\\
F_3(x,y)=f(x,y)|g(x,y)|,\quad F_4(x,y)=|f(x,y)|g(x,y).
\end{gathered}
\]
说法正确的是
\begin{choices}[columns=1]
\item 只有 $F_1(x,y)$ 在点 $P_0$ 处可微．
\item 只有 $F_1(x,y),F_2(x,y)$ 在点 $P_0$ 处可微．
\item 只有 $F_1(x,y),F_2(x,y),F_3(x,y)$ 在点 $P_0$ 处可微．
\item 4个函数均在点 $P_0$ 处可微．
\end{choices}
\end{question}
\QuestionBreak

\begin{question}[index=3]
已知幂级数 $\displaystyle\sum_{n=0}^{\infty}\frac{(x-a)^n}{\ln(n+2)}$
在点 $x_1=-2$ 处条件收敛，则 $\displaystyle\sum_{n=0}^{\infty}\frac{(x-a)^n}{(n+2)^2}$ 在点 $x_2=\dfrac12$ 处
\begin{choices}[columns=2]
\item 发散．
\item 绝对收敛．
\item 条件收敛．
\item 无法确定敛散性．
\end{choices}
\end{question}
\QuestionBreak

\begin{question}[index=4]
设 $\Omega$ 是锥体 $z\geq\sqrt{x^2+y^2}$ 介于 $z=1$ 和 $z=2$ 之间的部分，则
$\displaystyle\iiint_{\Omega} f(x^2+y^2+z^2)\dd v=$
\begin{choices}[columns=2,label-pos=left]
\item $\displaystyle\int_1^2\dd z\int_0^{2\pi}\dd\theta\int_0^z f(r^2+z^2)r\dd r$．
\item $\displaystyle\int_0^{2\pi}\dd\theta\int_1^2 r\dd r\int_0^1 f(r^2+z^2)\dd z$．
\item $\displaystyle\int_0^{2\pi}\dd\theta\int_0^{\frac\pi4}\dd\varphi\int_1^2 f(r^2)r^2\sin\varphi\dd r$．
\item $\displaystyle\int_0^{2\pi}\dd\theta\int_{\frac\pi4}^{\frac\pi2}\dd\varphi\int_1^2 f(r^2)r^2\sin\varphi\dd r$．
\end{choices}
\end{question}
\QuestionBreak

\begin{question}[index=5]
设 $\symbfit A,\symbfit B$ 均为3阶可逆实矩阵，$\symbfit A^*,\symbfit B^*$ 分别为 $\symbfit A,\symbfit B$ 的伴随矩阵，则
$\begin{pmatrix}\symbfit O&\symbfit A\\\symbfit B&\symbfit C\end{pmatrix}$ 的伴随矩阵为
\begin{choices}[columns=2,label-pos=left]
\item $\begin{pmatrix}-\symbfit B^*\symbfit C\symbfit A^*&-|\symbfit A|\symbfit B^*\\|\symbfit B|\symbfit A^*&\symbfit O\end{pmatrix}$．
\item $\begin{pmatrix}\symbfit B^*\symbfit C\symbfit A^*&-|\symbfit A|\symbfit B^*\\-|\symbfit B|\symbfit A^*&\symbfit O\end{pmatrix}$．
\item $\begin{pmatrix}-\symbfit A^*\symbfit C\symbfit B^*&|\symbfit A|\symbfit B^*\\|\symbfit B|\symbfit A^*&\symbfit O\end{pmatrix}$．
\item $\begin{pmatrix}\symbfit A^*\symbfit C\symbfit B^*&|\symbfit A|\symbfit B^*\\-|\symbfit B|\symbfit A^*&\symbfit O\end{pmatrix}$．
\end{choices}
\end{question}
\QuestionBreak

\begin{question}[index=6]
已知二次型 $f(x_1,x_2,x_3)$ 在正交变换 $\symbfit x=\symbfit Q\symbfit y$ 下化为 $y_1^2-y_2^2$，其中 $\symbfit Q=(\symbfit e_1,\symbfit e_2,\symbfit e_3)$．若
$\symbfit P=(\symbfit e_2,-\symbfit e_3,\symbfit e_1)$，则 $f(x_1,x_2,x_3)$ 在变换 $\symbfit x=\symbfit P\symbfit z$ 下化为
\begin{choices}[columns=4]
\item $z_1^2+z_2^2+z_3^2$．
\item $z_1^2-z_2^2+z_3^2$．
\item $-z_1^2+z_3^2$．
\item $-z_1^2-z_3^2$．
\end{choices}
\end{question}
\QuestionBreak

\begin{question}[index=7]
设 $\symbfit A$ 为 $m$ 阶正定矩阵，$\symbfit B$ 为 $m\times n$ 实矩阵，$\symbfit C=\symbfit B^{\mathrm T}\symbfit A\symbfit B$，则 $\symbfit C$ 与 $n$ 阶单位矩阵 $\symbfit E$ 合同的充分必要条件为
\begin{choices}[columns=1]
\item 齐次线性方程组 $\symbfit B^{\mathrm T}\symbfit x=\symbf{0}$ 有非零解．
\item 齐次线性方程组 $\symbfit B\symbfit x=\symbf{0}$ 只有零解．
\item 齐次线性方程组 $\symbfit B\symbfit B^{\mathrm T}\symbfit x=\symbf{0}$ 只有零解．
\item 齐次线性方程组 $\symbfit B^{\mathrm T}\symbfit B\symbfit x=\symbf{0}$ 有非零解．
\end{choices}
\end{question}
\QuestionBreak

\begin{question}[index=8]
设 $A_1,A_2,B$ 为随机事件，$0<P(B)<1$，$P(B)=P(A_1)P(B\mid A_1)+P(A_2)P(B\mid A_2)$，则以下结论：\par
① $A_1,A_2$ 是完备事件组；\par
② $P(A_1\mid B)+P(A_2\mid B)=1$；\par
③ $P(A_1\mid\overline B)+P(A_2\mid\overline B)=1$；\par
④ $P(A_1\mid B)+P(A_2\mid B)=P((A_1\cup A_2)\mid B)$．\par
正确结论的个数为
\begin{choices}[columns=4]
\item $1$．
\item $2$．
\item $3$．
\item $4$．
\end{choices}
\end{question}
\QuestionBreak

\begin{question}[index=9]
设随机试验 $E$ 有三种两两不相容的结果 $A,B,C$，且每次试验三种结果发生的概率分别为 $\dfrac16,\dfrac13,\dfrac12$．将试验 $E$ 独立重复做2次，$X$ 表示2次试验中结果 $A$ 发生的次数，$Y$ 表示2次试验中结果 $B$ 发生的次数，则 $X$ 与 $Y$ 的相关系数为
\begin{choices}[columns=4,label-pos=left]
\item $-\dfrac15$．
\item $-\dfrac{\sqrt5}{5}$．
\item $-\dfrac{\sqrt{10}}{10}$．
\item $-\dfrac{1}{10}$．
\end{choices}
\end{question}
\QuestionBreak

\begin{question}[index=10]
设总体 $X$ 服从参数为 $\lambda(\lambda>0)$ 的泊松分布，$X_1,X_2,\ldots,X_n$ 是来自总体 $X$ 的简单随机样本，则统计量
$\displaystyle T_1=\frac1n\sum_{i=1}^{n}X_i$，
$\displaystyle T_2=\frac1{n-1}\sum_{i=1}^{n-1}X_i+\frac1n X_n$ 满足
\begin{choices}[columns=2]
\item $ET_1>ET_2,\ DT_1>DT_2$．
\item $ET_1>ET_2,\ DT_1<DT_2$．
\item $ET_1<ET_2,\ DT_1>DT_2$．
\item $ET_1<ET_2,\ DT_1<DT_2$．
\end{choices}
\end{question}
